\chapter{Horizontal Velocity Interpolation}\label{ch:hvel}

Once a particle has been located (\cref{ch:localization}), the horizontal velocity must be
interpolated from the surrounding mesh values to the particle position. \textsf{KELP} provides
three schemes, selected by the configuration key \conf{uv\_interp}: a \emph{barycentric} scheme,
which blends the values at the vertices of the containing element; a \emph{least-squares} scheme
(the default), which fits a plane to the velocities at the element centres around the nearest node;
and an \emph{edge least-squares} scheme, which fits the plane to the containing element and its
edge neighbours, with a wall condition at the coast. The schemes are matched to the storage layout
of the velocity (\cref{sec:mesh-anatomy,tab:hvel-layouts}): the two least-squares schemes read
element-centre values, so supplying node-layout velocity to either is refused at startup. This
chapter gives the three schemes explicitly, including the local planar frames in which they compute,
and states the accuracy each achieves.

\begin{table}[t]
\centering\footnotesize\setlength{\tabcolsep}{4pt}
\begin{tabular}{@{}>{\raggedright\arraybackslash}p{58mm}>{\raggedright\arraybackslash}p{16mm}>{\raggedright\arraybackslash}p{50mm}@{}}
\toprule
velocity source & layout & admissible \conf{uv\_interp} \\
\midrule
FESOM2 \code{u}/\code{v}; FESOM-C element fields; HEALPix pixel geometry & element & \conf{leastsq}, \conf{edge\_leastsq}, \conf{barycentric} \\
FESOM2 \code{unod}/\code{vnod}; FESOM-C node fields; HEALPix dual mesh; NEMO (C-grid faces averaged to T-points) & node & \conf{barycentric} \\
\bottomrule
\end{tabular}
\caption{Velocity layouts by source and the interpolation schemes each admits. The layout is detected
from the field's dimensions; a metadata-only probe before the output file is created and a backstop
on the first step reject a mismatch with a message naming the fix.}
\label{tab:hvel-layouts}
\end{table}

\section{Projection to a local plane}\label{sec:hvel-projection}

All three schemes compute in a local planar frame rather than directly in longitude and latitude, so
that the spherical geometry is treated correctly away from the equator and across the dateline. The
two families use different frames.

\paragraph{Least-squares family.} The frame is centred on the particle $(\lonx_p,\latx_p)$ and scaled
by the metric at the particle's latitude,
\begin{equation}\label{eq:hvel-projection}
  z_0 = \cos\latx_p\cdot M, \qquad
  x_k = (\lonx_k-\lonx_p)\,z_0, \qquad
  y_k = (\latx_k-\latx_p)\,M ,
\end{equation}
with $M=\SI{111120}{\metre}$ per degree of latitude. The factor $\cos\latx_p$ is the essential
spherical correction: it shrinks the longitudinal axis so that a degree of longitude and a degree of
latitude map to their physical lengths. Across the dateline the plain least-squares scheme shifts an
element centre by one full turn only when the raw difference $\lonx_k-\lonx_p$ exceeds $\pm180^\circ$,
which leaves every stencil away from the seam bit-identical to the original arithmetic; the edge
least-squares scheme wraps every difference to $[-180^\circ,180^\circ)$.

\paragraph{Barycentric scheme.} The frame is centred on the element: its centre is the mean of the
vertices, with the longitudes accumulated as wrapped differences from the first vertex so that an
element straddling the dateline keeps its shape, and a point maps to
\begin{equation}\label{eq:hvel-projection-bary}
  x_k = \Delta\lonx_k \cdot \frac{\pi}{180}\, R \cos\latx_0, \qquad
  y_k = (\latx_k-\latx_0)\cdot \frac{\pi}{180}\, R ,
\end{equation}
where $(\lonx_0,\latx_0)$ is the frame centre, $\Delta\lonx_k$ the wrapped longitude difference and
$R$ the mesh's Earth radius. The weights computed below are ratios of areas in this plane and are
therefore independent of $R$; the radius serves only to keep the coordinates dimensionally
consistent.

\codeanchor{The projections are performed inline in \fn{\_interp\_leastsq\_core},
\fn{\_interp\_edge\_leastsq\_core} and \fn{\_interp\_barycentric\_core}
(\srcfile{numerics/interpolate.py}).}

\section{Barycentric interpolation}\label{sec:hvel-bary}

The barycentric scheme is the one scheme that serves node-layout velocity. The value at the particle
is the weighted average of the velocities at the vertices of the containing element, with weights
equal to the barycentric coordinates of the particle within it (\cref{fig:bary}); we give the
triangle here and the quadrilateral in \cref{sec:hvel-quad}. Let the projected triangle have vertices $1,2,3$ at
$(x_i,y_i)$ carrying velocities $(u_i,v_i)$, and let the particle project to
$(x_p,y_p)$. The barycentric coordinate $\bwt{i}$ is the area of the sub-triangle
\emph{opposite} vertex $i$ --- the sub-triangle formed by the particle and the other
two vertices --- divided by the area of the whole triangle:
\begin{equation}\label{eq:bary-area}
  \bwt{i} = \frac{\operatorname{area}(\text{sub-triangle opposite } i)}
                 {\operatorname{area}(\text{triangle})}, \qquad
  \bwt{1}+\bwt{2}+\bwt{3}=1 .
\end{equation}
A particle coinciding with vertex $i$ therefore has $\bwt{i}=1$ and the other weights
zero, and a particle on an edge has zero weight for the opposite vertex; the weights
vary continuously across element boundaries. In implemented form the weights are the
signed-area ratios
\begin{align}\label{eq:bary-explicit}
  D    &= (y_2-y_3)(x_1-x_3)+(x_3-x_2)(y_1-y_3), \notag\\
  \bwt{1} &= \frac{(y_2-y_3)(x_p-x_3)+(x_3-x_2)(y_p-y_3)}{D}, \\
  \bwt{2} &= \frac{(y_3-y_1)(x_p-x_3)+(x_1-x_3)(y_p-y_3)}{D}, \qquad
  \bwt{3} = 1-\bwt{1}-\bwt{2}, \notag
\end{align}
in which $D$ is twice the signed area of the triangle; a triangle whose $|D|$ falls
below a small tolerance is treated as degenerate and the interpolation reports
failure rather than dividing by it.

\begin{figure}[t]
\centering
\begin{tikzpicture}[scale=1.35]
  \coordinate (N1) at (0.0,0.0);
  \coordinate (N2) at (4.0,0.6);
  \coordinate (N3) at (1.6,3.1);
  \coordinate (P)  at (1.95,1.25);
  % sub-triangles, one shade each
  \fill[lagblue!12]  (P)--(N2)--(N3)--cycle;   % opposite N1 -> w1
  \fill[lagelem!16]  (P)--(N3)--(N1)--cycle;   % opposite N2 -> w2
  \fill[laggreen!14] (P)--(N1)--(N2)--cycle;   % opposite N3 -> w3
  \draw[meshedge,line width=1pt] (N1)--(N2)--(N3)--cycle;
  \draw[densely dotted,gray!70] (P)--(N1) (P)--(N2) (P)--(N3);
  % sub-area labels
  \node[font=\small,lagblue!60!black]  at (2.45,1.75) {$A_1\!\propto\!\bwt{1}$};
  \node[font=\small,lagelem!65!black]  at (1.05,1.35) {$A_2\!\propto\!\bwt{2}$};
  \node[font=\small,laggreen!50!black] at (2.0,0.55)  {$A_3\!\propto\!\bwt{3}$};
  % velocity arrows at nodes
  \draw[velarrow] (N1)--++(0.55,0.22); \draw[velarrow] (N2)--++(0.5,0.35);
  \draw[velarrow] (N3)--++(0.45,0.3);
  % nodes & drifter
  \node[meshnode] at (N1) {}; \node[meshnode] at (N2) {}; \node[meshnode] at (N3) {};
  \node[drifter] at (P) {};
  \node[nodelbl,below left=-1pt of N1]  {$n_1:(u_1,v_1)$};
  \node[nodelbl,right=1pt of N2]        {$n_2:(u_2,v_2)$};
  \node[nodelbl,above=0pt of N3]        {$n_3:(u_3,v_3)$};
  \node[font=\footnotesize,kelp!70!black,below right=-1pt of P] {drifter $P$};
\end{tikzpicture}
\caption{Barycentric interpolation of node-layout velocity. The three triangle
vertices $n_1,n_2,n_3$ carry the node velocities $(u_i,v_i)$. The drifter $P$
splits the triangle into three sub-triangles; the weight $\bwt{i}$ is the area of
the sub-triangle \emph{opposite} vertex $i$ divided by the total area, so the
nearer the particle lies to a vertex, the larger that vertex's weight. The
interpolated velocity is $\bwt{1}(u_1,v_1)+\bwt{2}(u_2,v_2)+\bwt{3}(u_3,v_3)$.}
\label{fig:bary}
\end{figure}

The interpolated velocity is then the weighted sum of the three vertex velocities,
multiplied by the integration-direction factor of \cref{sec:arch-direction}:
\begin{equation}\label{eq:bary-blend}
  (U,V) = \tdir \sum_{i=1}^{3} \bwt{i}\,(u_i,v_i).
\end{equation}
For node-layout fields the vertex velocity $(u_i,v_i)$ is read directly from the
node. For element-layout fields handled through the same routine, each vertex value
is first formed as the area-weighted average of the velocities of the elements
sharing that vertex, $u_n=\sum_{e\ni n}A_e u_e\,/\sum_{e\ni n}A_e$ (a vertex whose areas sum to zero
fails the interpolation), and the blend \cref{eq:bary-blend} is applied to those vertex values. The scheme reproduces a
constant field exactly in both layouts and a linearly varying field exactly in the
node layout; in the element layout the vertex-averaging step makes it
smoothing rather than linearly exact. The interpolated field is continuous across element edges.
A quadrilateral element is handled by the same blend with Wachspress coordinates
\citep{warren1996barycentric,floater2003mvc} in place of the triangular weights of
\cref{eq:bary-area} (\cref{sec:hvel-quad}).

\codeanchor{The barycentric blend is \fn{\_interp\_barycentric\_core}
(\srcfile{numerics/interpolate.py}); the triangle weights
\cref{eq:bary-explicit} are \fn{\_barycentric\_weights\_tri} and the Wachspress
generalization \fn{\_wachspress\_weights\_quad} (\srcfile{core/mesh.py}).}

\section{Least-squares interpolation (element layout)}\label{sec:hvel-leastsq}

For element-layout velocity, the default scheme fits a plane to the velocities of
the element centres surrounding the nearest node and evaluates that plane at the
particle (\cref{fig:leastsq}). The stencil is the one-ring of elements about the
nearest node; their centres, carrying the element velocities $(u_e,v_e)$, are the
data. With the plane written $f(x,y)=a_0+a_1 x + a_2 y$ in the particle-centred
frame of \cref{eq:hvel-projection} (so the particle sits at the origin), the
coefficients minimize the residual
\begin{equation}\label{eq:leastsq}
  \min_{a_0,a_1,a_2}\ \sum_{e} \bigl(a_0 + a_1 x_e + a_2 y_e - f_e\bigr)^2 ,
\end{equation}
and because the frame is centred on the particle the interpolated value is simply
the constant term, $U=\tdir\,a_0^{(u)}$ and $V=\tdir\,a_0^{(v)}$. The minimization
reduces to a symmetric $3\times3$ system of normal equations, solved directly by
Cramer's rule. Centring the coordinates on the particle is not cosmetic: it avoids the
catastrophic cancellation that would arise from forming the normal equations with
raw longitude and latitude values. The system is declared singular, and the interpolation fails, when
\begin{equation}\label{eq:leastsq-guard}
  |\det A|<10^{-30}\quad\text{or}\quad |\det A|<10^{-12}\,|a_{11}a_{22}a_{33}| ,
\end{equation}
the second test being relative to the scale of the metric entries: a stencil of fewer than three
non-collinear centres --- a coastal node with only two elements, for instance --- has a determinant
of pure rounding noise, which an absolute test alone would pass. The same guard protects the
plane fits of the vertical velocity, the column height and the diffusivity; the scalar-sampling fit uses
an absolute threshold of $10^{-12}$ with the same relative test. The scheme is
exact for fields varying linearly over the stencil. Its stencil is the one-ring of the \emph{nearest
node}, so the interpolated field jumps where the nearest node changes, and the anchor must be exact
for two runs to agree (\cref{sec:loc-exact}).

\begin{figure}[t]
\centering
\begin{tikzpicture}[scale=1.2]
  \coordinate (C)  at (2.00,1.55);
  \coordinate (p1) at (0.35,1.50);
  \coordinate (p2) at (1.15,0.20);
  \coordinate (p3) at (2.95,0.35);
  \coordinate (p4) at (3.70,1.70);
  \coordinate (p5) at (2.75,2.95);
  \coordinate (p6) at (1.00,2.80);
  \foreach \a/\b in {p1/p2,p2/p3,p3/p4,p4/p5,p5/p6,p6/p1}{
     \fill[meshface] (C)--(\a)--(\b)--cycle;
     \draw[meshedge] (C)--(\a); \draw[meshedge] (\a)--(\b);}
  % element centres of the one-ring (the leastsq data)
  \foreach \a/\b in {p1/p2,p2/p3,p3/p4,p4/p5,p5/p6,p6/p1}
     \node[elemcentre] at ($ 0.3333*(C)+0.3333*(\a)+0.3333*(\b) $) {};
  \foreach \p in {C,p1,p2,p3,p4,p5,p6} \node[meshnode] at (\p) {};
  \node[drifter] (D) at (2.25,1.5) {};
  \node[nodelbl,above left=0pt of C] {nearest node};
  \node[elemlbl] at (3.15,1.05) {element centres $(u_e,v_e)$};
  \node[font=\footnotesize,kelp!70!black,below=1pt of D] {drifter};
\end{tikzpicture}
\caption{Least-squares interpolation of element-layout velocity. The data are the
velocities at the centres (orange squares) of the one-ring of elements about the
nearest node. A plane $f=a_0+a_1x+a_2y$ is fitted to them in a frame centred on the
drifter, so the interpolated velocity is the fitted constant $a_0$. The scheme is
exact for linearly varying fields over this stencil.}
\label{fig:leastsq}
\end{figure}

The schemes are tied to the storage layout, and the engine enforces the match rather than silently
producing a wrong answer. Because both least-squares schemes index velocities at element centres,
supplying them with node-layout velocities such as \conf{unod}/\conf{vnod} is a category error; the
engine detects the mismatch from the field's dimensions and raises a configuration error at startup,
with the message that node-based velocities require \conf{uv\_interp}${}={}$\conf{barycentric}. The
check is performed both during a metadata-only probe before any data is read and as a backstop on the
first step, so the run fails fast and informatively rather than midway through.

\codeanchor{The plane fit is \fn{\_interp\_leastsq\_core}
(\srcfile{numerics/interpolate.py}); the layout-compatibility check that rejects
node velocities under both least-squares schemes is \fn{\_validate\_uv\_compat}
(\srcfile{tracking.py}), fed by each reader's \fn{probe\_uv\_layout}.}

\section{Edge least-squares interpolation}\label{sec:hvel-edge}

The edge least-squares scheme (\conf{uv\_interp: edge\_leastsq}) fits the same plane as
\cref{eq:leastsq}, in the particle-centred frame, but over a different stencil: the centre of the
element that contains the particle and the centres of the elements across each of its edges --- four
points for a triangle, five for a quadrilateral (\cref{fig:edge-lsq}). The stencil follows the
containing element rather than the nearest node, as in the finite-volume discretisation of the FESOM
models, and it needs the containing element, so a particle outside every element fails the
interpolation. The plane is exact for a linear field in the interior.

At the coast an edge has no neighbour. It contributes a \emph{ghost} point instead: the own centre
$\mathbf{c}$ mirrored across the edge, $\mathbf{g}=2\mathbf{f}-\mathbf{c}$, with $\mathbf{f}$ the foot of
$\mathbf{c}$ on the edge. With $\hat{\mathbf{t}}$ the unit tangent of the edge and $\mathbf{u}$ the own
velocity, the ghost carries
\begin{equation}\label{eq:edge-ghost}
  \mathbf{u}_g =
  \begin{cases}
    2(\mathbf{u}\cdot\hat{\mathbf{t}})\,\hat{\mathbf{t}}-\mathbf{u} & \text{free slip (default)},\\
    -\mathbf{u} & \text{no slip},
  \end{cases}
\end{equation}
selected by \conf{edge\_leastsq\_bc}. The free-slip ghost is the mirror image of the own velocity ---
its tangential component kept and its normal component reversed --- so the fitted plane carries a
flow parallel to the wall exactly and damps the flow through it. The no-slip ghost makes the plane
vanish near the wall; being a single plane fitted to both sides, it overshoots on the interior side of
a wall element and can reverse the flow near a corner between two walls, a property of the fit rather
than a defect. A boundary edge of zero length contributes no point. The stencil stays well conditioned
even on very thin elements, because the mirrored ghosts span both directions.

\begin{figure}[t]
\centering
\begin{tikzpicture}[scale=1.25]
  \coordinate (A) at (0,0); \coordinate (B) at (3,0); \coordinate (C) at (1.3,2.2);
  \coordinate (D) at (3.6,2.3); \coordinate (E) at (-1.2,1.8);
  \fill[pattern=north east lines,pattern color=gray!45] (-1.4,0)--(4.0,0)--(4.0,-1.25)--(-1.4,-1.25)--cycle;
  \fill[meshface] (B)--(C)--(D)--cycle; \fill[meshface] (C)--(A)--(E)--cycle;
  \fill[hilface] (A)--(B)--(C)--cycle;
  \draw[meshedge] (B)--(D)--(C) (C)--(E)--(A);
  \draw[laggreen,line width=1.2pt] (-1.4,0)--(4.0,0);
  \node[font=\footnotesize,laggreen!70!black,fill=white,inner sep=1pt,anchor=north west] at (-1.35,-0.08) {coast (boundary edge)};
  \foreach \p in {A,B,C,D,E} \node[meshnode] at (\p) {};
  \node[elemcentre] (c) at (1.433,0.733) {};
  \node[elemcentre] at (2.633,1.5) {};
  \node[elemcentre] at (0.033,1.333) {};
  \node[elemcentre,fill=white] (g) at (1.433,-0.733) {};
  \draw[densely dotted,gray] (c)--(g);
  \node[circle,fill=black,inner sep=0.8pt] at (1.433,0) {};
  \node[font=\footnotesize,anchor=south west] at (1.46,0.0) {$\mathbf{f}$};
  \draw[velarrow] (c)--++(0.55,-0.28);
  \draw[velarrow] (g)--++(0.55,0.28);
  \node[font=\footnotesize,anchor=east] at (1.36,0.78) {$\mathbf{c}$};
  \node[font=\footnotesize,fill=white,inner sep=1pt,anchor=east] at (1.3,-0.78) {$\mathbf{g}=2\mathbf{f}-\mathbf{c}$};
  \node[drifter] at (1.75,1.05) {};
  \node[font=\footnotesize,fill=white,inner sep=1pt,anchor=west] at (2.1,-0.55) {$\mathbf{u}_g=2(\mathbf{u}\cdot\hat{\mathbf{t}})\hat{\mathbf{t}}-\mathbf{u}$};
  \node[elemlbl,anchor=west] at (2.75,1.55) {neighbour centres};
\end{tikzpicture}
\caption{The edge least-squares stencil. The drifter lies in the shaded triangle; the plane is fitted
to its centre $\mathbf{c}$ and the centres of the elements across its two interior edges. Its third
edge is on the coast and contributes a ghost at the mirror image $\mathbf{g}$ of $\mathbf{c}$, carrying
the mirrored velocity: the component along the wall is kept and the component into the wall reversed
(free slip), or the whole velocity reversed (no slip).}
\label{fig:edge-lsq}
\end{figure}

\codeanchor{\fn{\_interp\_edge\_leastsq\_core} (\srcfile{numerics/interpolate.py}); the edge
neighbours are \fn{build\_elem\_neighbors} (\srcfile{core/mesh.py}), built once when the run
starts.}

\section{Interpolation on mixed meshes: quadrilateral elements}\label{sec:hvel-quad}

The barycentric scheme of \cref{sec:hvel-bary} generalizes from the triangular meshes of FESOM2 to
meshes with quadrilaterals: the mixed meshes of FESOM-C, the dual mesh built from a HEALPix grid
(mostly quadrilaterals; an all-wet sphere has exactly eight triangles), the pixel geometry of a
HEALPix grid and the meshes built from structured NEMO grids (quadrilaterals only). Every horizontal
interpolator is split into a \emph{geometric weight stage}, which depends on the number of vertices of
the containing element, and a \emph{value stage}, which depends only on the field layout. The plain
least-squares scheme bypasses the geometric stage entirely --- it reads one velocity per element
centre --- and is identical on triangles and quadrilaterals; the edge least-squares scheme differs only
in the number of edges.

Before the weight stage, the element's vertices are compacted: the $-1$ padding is removed, and a
triangle that FESOM-C has stored as a quadrilateral with a repeated leading vertex is reduced to its
three distinct vertices. A triangle on a mixed mesh is therefore interpolated by the triangular
weights of \cref{eq:bary-explicit}, bit-for-bit as on a pure FESOM2 mesh; only a genuine four-vertex
element takes the quadrilateral path.

For a quadrilateral the weights are the Wachspress generalized barycentric coordinates
\citep{warren1996barycentric,floater2003mvc}, the natural analogue of \cref{eq:bary-area} for a convex
polygon (\cref{fig:wachspress}). With the four vertices $u_1,\dots,u_4$ in cyclic order, the particle
$p$, and the doubled signed areas
\begin{equation}\label{eq:wachspress}
  A_k = 2\operatorname{area}(p,u_k,u_{k+1}), \qquad
  C_k = 2\operatorname{area}(u_{k-1},u_k,u_{k+1}),
\end{equation}
the classical weight is $C_k/(A_{k-1}A_k)$, normalized over the vertices. The engine uses its
\emph{product form}, the classical weight multiplied by $A_1A_2A_3A_4$,
\begin{equation}\label{eq:wachspress-product}
  \bwt{k} = \frac{C_k\,A_{k+1}\,A_{k+2}}{\sum_{j} C_j\,A_{j+1}\,A_{j+2}} ,
\end{equation}
which contains no division by an area that vanishes on an edge: on an edge it reduces exactly to
linear interpolation along that edge and at a vertex to that vertex's indicator. Each term carries
three signed areas, so a clockwise element negates every term and the normalization cancels it; no
orientation handling is needed. The coordinates reproduce a linear field exactly, are smooth in the
interior and give every vertex a non-zero weight.

\begin{figure}[t]
\centering
\begin{tikzpicture}[scale=1.25]
  \coordinate (v1) at (0.2,0.2);
  \coordinate (v2) at (3.6,0.0);
  \coordinate (v3) at (3.2,2.9);
  \coordinate (v4) at (0.0,2.6);
  \coordinate (P)  at (1.70,1.35);
  \fill[lagblue!12]  (P)--(v1)--(v2)--cycle;
  \fill[lagelem!16]  (P)--(v2)--(v3)--cycle;
  \fill[laggreen!14] (P)--(v3)--(v4)--cycle;
  \fill[gray!15]     (P)--(v4)--(v1)--cycle;
  \draw[meshedge,line width=1pt] (v1)--(v2)--(v3)--(v4)--cycle;
  \draw[densely dotted,gray!70] (P)--(v1) (P)--(v2) (P)--(v3) (P)--(v4);
  \node[font=\small] at (1.83,0.50) {$A_1$};
  \node[font=\small] at (2.78,1.42) {$A_2$};
  \node[font=\small] at (1.63,2.20) {$A_3$};
  \node[font=\small] at (0.70,1.38) {$A_4$};
  \node[meshnode] at (v1) {}; \node[meshnode] at (v2) {};
  \node[meshnode] at (v3) {}; \node[meshnode] at (v4) {};
  \node[drifter] at (P) {};
  \node[nodelbl,below left=-1pt of v1]  {$u_1$};
  \node[nodelbl,below right=-1pt of v2] {$u_2$};
  \node[nodelbl,above right=-1pt of v3] {$u_3$};
  \node[nodelbl,above left=-1pt of v4]  {$u_4$};
  \node[font=\footnotesize,kelp!70!black,right=1pt of P] {$p$};
\end{tikzpicture}
\caption{Wachspress interpolation on a quadrilateral element. The drifter $p$ splits the element into
four fan triangles of doubled area $A_1,\dots,A_4$. The weight at vertex $k$ is proportional to
$C_k A_{k+1} A_{k+2}$, where $C_k$ is the doubled area of the corner triangle $(u_{k-1},u_k,u_{k+1})$:
the two fan triangles \emph{not} adjacent to vertex $k$ grow as $p$ approaches it.}
\label{fig:wachspress}
\end{figure}

The formula needs no guard and has no fall-back because quadrilaterals are required to be convex:
the FESOM-C reader, the builder of structured-grid meshes and the HEALPix builders check every
quadrilateral when the mesh is loaded and name the offending elements if one is not convex.

The \emph{containment} test of \cref{sec:loc-contain}, which decides whether a particle lies inside a
quadrilateral, uses a fan of triangles --- never the Wachspress construction, which enters only the
weight stage after containment has been established. The same vertex-count branch governs the
\conf{barycentric} option of the vertical-velocity combine (\cref{sec:vert-whoriz}) and of scalar
sampling (\cref{ch:sampling}); the \conf{nearest}, \conf{idw} and \conf{leastsq} options of those
routines are independent of element shape. Element areas enter only when an element-layout vertex
value is formed as an area-weighted average (\cref{sec:hvel-bary}). They are read from the model's mesh
file when it provides them (FESOM2 mesh diagnostics do; ASCII meshes and regional extracts do not); otherwise they are computed on the sphere as the
spherical excess of the element's fan of triangles times the square of the mesh radius. Because they
appear only as ratios, their absolute scale is immaterial.

\section{Common properties}\label{sec:hvel-common}

\paragraph{Three dimensions.} In a three-dimensional run the horizontal velocity is piecewise constant
in the vertical: each particle reads the layer found in \cref{sec:vert-layer}, and the scheme is applied
to that layer's field without vertical interpolation.

\paragraph{Time.} Between two snapshots the time-aware variants blend the two fields in time for each
element or vertex they read, before the unchanged spatial formula (\cref{ch:temporal,ch:fused}).

\paragraph{Failure.} A scheme that cannot interpolate --- a singular fit, no containing element, a
degenerate triangle --- reports failure; the particle's step is undone, its status set to
\code{OUTSIDE} and its failure reason to an interpolation failure (\cref{ch:lifecycle}).

\paragraph{Element tracking.} With \conf{elem\_search: walk} (\cref{sec:loc-walk}) the barycentric and
edge least-squares schemes receive the containing element from the walk and skip their own search; the
plain least-squares scheme does not use it.

\codeanchor{The weight dispatcher is \fn{\_polygon\_weights}; the triangular and Wachspress weights are
\fn{\_barycentric\_weights\_tri} and \fn{\_wachspress\_weights\_quad}; the convexity check is
\fn{validate\_quad\_convexity} with \fn{find\_nonconvex\_quads}; element areas are
\fn{build\_elem\_areas}; all in \srcfile{core/mesh.py}. The vertex compaction is in
\fn{\_interp\_barycentric\_core} (\srcfile{numerics/interpolate.py}).}

\codenote{The projections use three radii: the barycentric horizontal-velocity kernel uses the mesh
radius, the containment test and the vertical-velocity and scalar barycentric kernels use
\SI{6371000}{\metre}, and the least-squares family uses \SI{111120}{\metre} per degree. The weights
are ratios of areas in the projected plane and the least-squares value is the fitted constant, so the
interpolated values do not depend on the radius; we record the difference only for fidelity to the
source.}
